\documentclass[11pt]{article}
\usepackage[T1]{fontenc}
\usepackage[utf8]{inputenc}
\usepackage{lmodern}
\usepackage{microtype}
\usepackage{booktabs}
\usepackage{graphicx}
\usepackage{hyperref}
\usepackage[margin=1in]{geometry}

\title{Evidence-Bounded Generation for Low-Resource Languages:\\
Human-Reviewed Grammar, Corpus Attestation, and LLM Proposals}
\author{Fabrício Ferraz Gerardi}
\date{}

\begin{document}
\maketitle

\begin{abstract}
This paper presents an evidence-bounded architecture for controlled sentence
generation in a low-resource language, using Boe-Bororo as a case study. The
architecture separates corpus observation, human linguistic analysis,
grammatical licensing, deterministic realization, and subsequent corpus
attestation. Corpus occurrence therefore supplies evidence but does not
automatically license productive generation. A frozen baseline evaluates five
lexemes across 28 reviewed structural cells: all 28 are generated, none of the
complete generated sentences is an exact corpus sentence, and 14 unique
generated predicate forms are independently attested as corpus tokens. We then
compare direct LLM surface generation with an evidence-bounded condition in
which the same model proposes only an abstract structure and reviewed grammar
realizes the surface form. Direct generation yields one exact surface match in
28 tasks; this is not treated as a grammaticality score, because human audit
distinguishes morphological well-formedness from construction, person, and
coding-frame adequacy. In the evidence-bounded condition, all 28 proposals
preserve the requested structural specification and all 28 are realized by the
grammar frozen before model outputs were inspected. The results support a
division of labor in which an LLM can propose structures without serving as the
source of morphological or constructional knowledge, while unresolved evidence
remains explicitly unresolved.
\end{abstract}

\section{Introduction}

Generative experiments with low-resource languages face a basic evidential
problem: occurrence in a corpus is not equivalent to a grammatical rule, while
absence from a small corpus is not evidence of ungrammaticality. The problem
becomes more acute when automatic induction or a large language model is allowed
to convert distributional regularities directly into productive analyses.
Surface plausibility can then obscure whether an output is licensed by
documented grammatical evidence, extrapolated from an incomplete paradigm, or
simply produced by analogy.

We investigate an alternative architecture in which documentary corpus evidence
and grammatical licensing remain formally distinct. The CorBo corpus supplies
observations. Human linguistic review establishes the morphological and
constructional facts that may license generation. A deterministic generator
realizes only reviewed combinations and retains provenance for the evidence and
rules used. Unresolved evidence remains an explicit computational state rather
than being completed automatically.

The study evaluates this architecture in two stages. First, a frozen controlled
baseline asks whether a small reviewed grammar can recombine licensed Bororo
structures without reducing generation to retrieval from the corpus. Second, we
compare two uses of the same LLM over a common frozen task inventory. In the
direct condition, the model produces Bororo surface output. In the
evidence-bounded condition, it produces only an abstract structural proposal;
surface morphology is then supplied by the reviewed deterministic grammar or
the request is blocked if the required license is absent.

This comparison requires a distinction that is often hidden by exact-match
evaluation. A model output may differ from the requested realization while
remaining a possible expression in another construction. We therefore separate
morphological well-formedness from construction match, S/A/O person match, and
coding-frame match. Corpus attestation is measured independently and never
promotes an analysis to grammatical status.

The paper addresses three empirical questions. First, can reviewed but partial
grammatical knowledge support controlled recombination beyond exact corpus
sentences? Second, what kinds of divergence arise when an LLM is asked to
produce Bororo morphology directly? Third, can the model preserve the requested
structural information when morphology is delegated to an evidence-bounded
realizer? These questions concern evidential control and task preservation, not
a claim of complete Bororo coverage or a single model-wide grammaticality
score.

The contribution is therefore methodological. The architecture makes
\emph{observation}, \emph{analysis}, \emph{licensing}, \emph{generation},
and \emph{attestation} separately inspectable, while the frozen comparison
tests a division of labor in which an LLM can serve as a proposal mechanism
without becoming the source of the grammatical generalizations against which
its outputs are evaluated.

\section{Boe-Bororo and the corpus setting}
\label{sec:bororo}

Bororo is a Bororoan language spoken in Mato Grosso, Brazil. The grammatical
description underlying the present system is based especially on data from
speakers in Meruri. Contemporary Bororo communities are sociolinguistically
heterogeneous: their historical trajectories differ, and speaker competence
varies by generation and locality. Consequently, variation in documentary
materials cannot be interpreted independently of locality, chronology, contact,
mobility, and language shift.

This point also affects the treatment of historical documentation. Sources
produced in different periods and contact settings are not interchangeable
samples of a single invariant linguistic system. Historical records and modern
fieldwork differ in locality, transcription practice, consultant background,
genre, and documentary purpose. The controlled-generation experiment therefore
does not collapse all documentary forms into a single automatically productive
grammar.

The present experiment uses a reviewed subset of CorBo represented in CoNLL-U.
The corpus functions as the observational layer of the system. Token forms,
lemmas, UPOS labels, dependency relations, sentence contexts, and translations
may be retrieved for human inspection, but corpus frequency or formal
resemblance does not promote a morphological analysis to reviewed status. The
trusted experimental subset is deliberately separated from larger working
collections whose annotation remains under development.

This separation is important for reproducibility. The generator consumes a
specified corpus snapshot, while grammatical licenses are stored independently
and require explicit human review. Thus a change in corpus annotation can alter
the observational evidence without silently changing the grammar used for
generation.

\section{Corpus evidence is not grammatical licensing}
\label{sec:evidence}

A corpus establishes that a form occurred in a documented context; a generator
requires a different decision, namely whether a reviewed analysis licenses that
form or construction for productive use. We therefore distinguish
\emph{observation}, \emph{analysis}, \emph{licensing}, \emph{generation},
and \emph{attestation}. Observation records corpus material, analysis records a
human linguistic interpretation, licensing authorizes a specific computational
operation, generation records the resulting candidate and provenance, and
attestation asks whether its surface or predicate form occurs independently in
the corpus. Attestation does not promote itself into a grammatical license.

The same separation applies to lexical identity and partial paradigms.
Homographic forms are not automatically assigned to one lexeme, and a
familiar-looking affix is not segmented solely from its shape. An exact reviewed
person cell may therefore coexist with an unresolved stem class, while a form
reviewed in an imperative, subjunctive, gerundial, negative, or irrealis
construction does not automatically establish the corresponding ordinary
indicative cell. These unresolved states are represented rather than completed
by analogy.

\section{Related work}
\label{sec:related}

Recent work on endangered-language NLP provides an important comparison for
the present architecture. Zhang et al.~\cite{zhang2022cherokee} frame
computational work on Cherokee around collaboration with language communities
and propose machine-in-the-loop methods for enriching resources. Joshi
et al.~\cite{joshi2019unsung} likewise emphasize that low-resource language
technology is not only a modeling problem: data creation, deployment, and the
conditions under which a technology is useful to a community are integral to
the task.

NüshuRescue provides a particularly close computational contrast
\cite{yang2025nushurescue}. Yang et al.\ construct a 500-sentence manually
validated Nüshu--Chinese corpus and use 35 seed examples with GPT-4-Turbo to
generate additional translations. Their pipeline combines LLM generation,
automatic validation, and manual review. A rule-based length validator enforces a one-to-one character mapping and
rejects length mismatches, prompting the model again for up to seven retries.
The authors report 31.37\% accuracy without enforced length validation and
48.69\% with it; the latter is measured as perfect character-by-character
translation of 50 held-out NCGold examples using 35 seed sentences. The study
therefore provides a concrete example in which explicit language-specific
validation improves an LLM-based low-resource pipeline.

The present work shares the broader human-in-the-loop motivation but addresses
a different problem. NüshuRescue uses an LLM primarily to expand parallel data
and seeks to reduce the amount of manual annotation required. Our system does
not treat generated data as a route to automatic grammar expansion. Instead,
human-reviewed linguistic analysis defines the space within which generation is
licensed. The LLM, when introduced, is positioned upstream as a proposal
mechanism; it cannot promote its own output into grammatical knowledge.

This contrast is consequential for evaluation. Translation accuracy measures
agreement with a reference translation. Controlled grammatical generation
requires additional questions: whether a person realization was independently
reviewed, whether a coding frame licenses the requested argument structure,
whether construction-specific morphology has been generalized beyond its
evidence, and whether corpus attestation is being mistaken for grammatical
productivity. Our evaluation therefore reports provenance and licensing
coverage separately from surface attestation.

\subsection{Computational documentation and explicit grammar}

A second line of work connects language documentation with reusable
computational descriptions. The CLD$^2$ proposal explicitly argues for a
two-way relation between documentation and NLP: documentary projects should
produce NLP-friendly annotated data, while computational work should exploit
documentary databases for endangered-language technology
\cite{zariquiey2022cld2}. The present architecture can be situated within this
interface, but adds an explicit review boundary between documentary evidence
and productive computational knowledge.

Finite-state work on Indigenous and endangered languages demonstrates the
value of linguistically explicit morphology under severe data scarcity. For
Yine, Ingunza Torres et al.~\cite{ingunza2021yine} encode morphophonology,
possession, and verbal predicates in finite-state transducers intended both for
documentation and downstream language technology. Comparable analyzers model
morphotactics and morphophonological alternations in Paraguayan Guaraní
\cite{kuznetsova2021guarani} and Highland Puebla Nahuatl
\cite{pugh2023nahuatl}. Work on Lakota is particularly relevant because a
finite-state morphological component is coupled to a precision HPSG grammar
and can operate in the generative direction \cite{curtis2014lakota}.

These systems establish that explicit linguistic descriptions remain useful
computational objects in low-resource settings. Our contribution is not a new
finite-state formalism. Instead, it concerns the evidential control surrounding
generation: lexical and constructional knowledge may remain partial, exact
reviewed cells coexist with fuller paradigms, corpus matches remain
observations, and unresolved evidence remains computationally unresolved.

AvarLab offers a recent and especially close point of comparison
\cite{zagidov2026avarlab}. Its rule-based generate--verify workflow generates
large inflectional inventories from lexical resources and couples morphological
generation with corpus verification. Our workflow shares the use of explicit
rules and corpus comparison. The present paper, however, asks a different
evidential question: corpus attestation is represented separately from the
human-reviewed license that permits productive generation. We therefore do not
treat a corpus match, by itself, as sufficient to promote a candidate into the
reviewed grammatical inventory. This separation is motivated by documentary
situations in which corpus sparsity, homography, construction-specific
allomorphy, and incomplete paradigms can make automatic generalization risky.

\subsection{Human-in-the-loop systems and the closest architectural comparison}

Human-in-the-loop NLP is usually formulated as a feedback cycle in which human
judgments improve a model or its training data \cite{wang2021hitl}. Recent
discussion of LLM-based applications emphasizes that this remains especially
important for low-resource languages, where traditional pipelines can provide
greater control and can operate with data regimes that do not support direct
LLM training \cite{eckert2026humans}. Our use of human review is narrower and
more formal: review changes an explicit inventory of computational licenses,
rather than implicitly updating model parameters.

A particularly close 2026 comparison is the \emph{langlit} platform
\cite{haberland2026langlit}. It combines a finite-state morphological analyzer,
a three-tier human-in-the-loop annotation workflow, searchable corpora,
grammar-guided interactive word construction, corpus-linked hypothesis
tracking with provenance, an editable dictionary, access controls, and optional
LLM integration. This substantially narrows the novelty claim available to the
present work: integration of corpus evidence, explicit morphology, provenance,
human review, and optional LLM assistance is not by itself novel.

The distinction investigated here is therefore deliberately narrower than
platform integration. The generator represents \emph{licensing status} as an
independent computational layer: corpus observation, human linguistic analysis,
a license for generation, a generated candidate, and subsequent corpus
attestation are distinct states. Partial knowledge is represented explicitly:
an exact reviewed person cell need not imply a paradigm, and a
construction-specific realization need not imply the corresponding ordinary
indicative cell. We evaluate this representation directly by freezing the
reviewed grammar before comparing direct LLM surface generation with abstract
LLM proposals whose realization is delegated to that grammar. The claimed
contribution is this operationalization and evaluation of evidence-bounded
licensing, rather than the individual presence of an FST or rule system,
corpus search, provenance, human review, or LLM assistance.

\section{Evidence-bounded architecture}
\label{sec:architecture}

The current pipeline is:

\begin{quote}
CorBo $\rightarrow$ corpus evidence $\rightarrow$ human linguistic review
$\rightarrow$ reviewed grammar $\rightarrow$ controlled generator
$\rightarrow$ generated candidate $\rightarrow$ validation and corpus
attestation.
\end{quote}

The arrows in this pipeline are deliberately asymmetric. Corpus evidence may
trigger human review, but only a review decision changes the inventory of
generation licenses. After generation, corpus attestation is reported as a
separate evidential property of the candidate; predicate-form attestation, for
example, does not establish that the complete generated structural cell was
observed.

\section{Reviewed grammatical constraints}
\label{sec:grammar}

The v1 grammar used in the experiments is intentionally narrower than the
descriptive grammar from which individual analyses are reviewed. It records
three coding-frame labels needed by the generator: \emph{monovalent},
\emph{extended-intransitive}, and \emph{divalent}. Coding-frame review is
independent of morphological review; similarity of surface form or translation
does not assign a frame.

Person morphology has two licensing scopes. A \texttt{full\_stem\_class}
license makes available the independently reviewed person-index paradigm for
that lexical assignment. An \texttt{exact\_person\_cell} license makes
available only the explicitly reviewed lemma--person realization. Exact cells
therefore do not complete missing members of a paradigm.

Construction-specific morphology is stored separately from ordinary indicative
person cells. A reviewed imperative, subjunctive, gerund, optative, negative, or
irrealis realization consequently licenses only that reviewed constructional
cell unless broader evidence has been entered independently. The morphotactic
component likewise separates ordering constraints from surface realization of
reviewed operator combinations. It records, for example, reviewed
IRR--IND and IRR--NEG--IND realizations without treating them as a general
fixed-slot template.

The v1 constraint inventory also records \textit{-re}, \textit{-wo},
\textit{-iagu}, \textit{-ie}, and \textit{-ia} as members of one reviewed
mutually exclusive suffix domain. This is a co-occurrence constraint only:
membership in the domain does not assign a functional analysis to a suffix
whose function has not independently been reviewed.

\section{Controlled generation}
\label{sec:generation}

The experiment uses the ordinary indicative dispatcher as the frozen
surface-realization path. The reviewed lexical frame determines which of three
generators is called; there is no fallback to another frame when the requested
combination lacks evidence. A missing required person cell, A host, or selected
oblique therefore returns a blocked result rather than an analogically completed
form.

For a reviewed monovalent predicate, the generator combines a licensed
S-indexed stem with indicative \textit{-re}. For a reviewed
extended-intransitive predicate, the same core predicate pattern is retained and
an additional participant may be added only through a lexically reviewed
postpositional relation. The existence of an extended-intransitive frame does
not itself license a particular postposition.

For a reviewed divalent predicate, the implementation keeps A and O as distinct
coding positions. The frozen v1 pattern is \texttt{(RP-A) A=re (RP-O) O=LEX}:
the A expression uses a separately reviewed indicative host, while the predicate
is realized with the reviewed O-indexed stem. Corpus occurrence of that
predicate form is therefore evidence about the predicate morphology, not
independent attestation of the complete A$\times$O structural cell.

This dispatcher is deliberately construction-specific. Other reviewed forms
present elsewhere in the computational grammar---for example imperative,
subjunctive, gerundial, optative, negative, or irrealis cells---do not enter the
v1 ordinary-indicative experiment merely because they are available to other
modules. The experimental surface inventory is thus fixed by the dispatcher and
its reviewed dependencies, not by every analysis represented in the broader
codebase.

\section{Experiment 1: controlled-generation baseline}
\label{sec:evaluation}

We froze a first experimental baseline after resolving the review queue for five
target lexemes: \textit{nudu}, \textit{meru}, \textit{kodu}, \textit{mako},
and \textit{maku}. A \emph{structural cell} in this report is a requested
combination of lexical item, reviewed coding frame, and the participant-person
values required by that frame. It is an evaluation unit, not a claim that a
complete sentence with those values has previously occurred in the corpus. The
28 requests were selected from combinations for which the frozen v1 dispatcher
had the licenses required for ordinary indicative realization.

\begin{table}[ht]
\centering
\caption{Controlled-generation baseline (v1; report schema 1.2).}
\label{tab:baseline}
\begin{tabular}{lr}
\toprule
Measure & Value\\
\midrule
Lexemes evaluated & 5\\
Structural cells requested & 28\\
Generated records & 28\\
Blocked requests & 0\\
Complete generated sentences exactly attested & 0\\
Unique generated predicate forms attested in corpus & 14\\
Records licensed by exact person cell & 20\\
Records licensed by full stem class & 8\\
Reviewed construction-specific cells in grammar & 16\\
\bottomrule
\end{tabular}
\end{table}

The inventory comprises 16 cells for three monovalent lexemes, six for one
extended-intransitive lexeme, and six for one divalent lexeme. All 28 requests
are realized and none is blocked. At the morphological licensing step, 20
records (71.4\%) use an exact reviewed lemma--person cell, whereas eight
(28.6\%) use an independently reviewed full stem-class assignment. These
figures describe the provenance of the generated records; they are not two
levels of grammatical confidence.

The report also counts 16 reviewed construction-specific cells in the frozen
grammar. This is an inventory statistic rather than a subset of the 28
ordinary-indicative experimental requests: those construction-specific cells
remain segregated from the dispatcher used for this baseline.

Sentence and predicate-form attestation are likewise distinct measurements.
Zero of the 28 complete generated outputs is an exact normalized sentence match
in the trusted corpus. Independently, 14 \emph{unique} generated predicate
forms occur as corpus tokens. The latter count is type-based and therefore
should not be read as 14 of 28 structural cells being attested. In particular,
predicate-form occurrence does not establish the complete participant
configuration of a generated cell, and for divalent predicates it does not
supply independent evidence for the generated A$\times$O combination.

The baseline consequently demonstrates controlled recombination under the
frozen licenses without exact sentence retrieval. Corpus matches remain
post-generation evidence: neither predicate-form presence nor sentence absence
changes the reviewed grammatical inventory used to produce the same records.

\section{Experiment 2: direct vs. evidence-bounded LLM generation}
\label{sec:ai}

The second experiment evaluates the central operational claim of the
evidence-bound architecture: probabilistic proposal can be separated from
grammatical realization so that model output does not itself supply the
morphological license. The comparison uses the same frozen lexical and
constructional inventory in two conditions.

\paragraph{Condition A: direct generation.}
An LLM receives a task specification and is asked to return a Bororo surface
form or sentence directly. The output is retained exactly as produced and is
subsequently audited against the reviewed grammar and corpus evidence.

\paragraph{Condition B: evidence-bounded generation.}
The same LLM receives an equivalent task but returns an abstract proposal
rather than Bororo surface morphology. A proposal specifies the lexeme,
construction, relevant S/A/O person values, and any oblique phrase represented
in the task schema. The deterministic component either realizes the proposal
from reviewed licenses or blocks it. The LLM is not permitted to add a person
cell, stem class, coding frame, morpheme analysis, or constructional realization
to the reviewed grammar.

\begin{table}[ht]
\centering
\caption{Comparison between direct and evidence-bounded generation.}
\label{tab:llmexperiment}
\begin{tabular}{p{.28\linewidth}p{.31\linewidth}p{.31\linewidth}}
\toprule
Property & Direct LLM & Evidence-bounded condition\\
\midrule
LLM produces Bororo surface & yes & no; proposes abstract request\\
Morphology source & model output & reviewed deterministic realization\\
Unsupported request & may surface & explicitly blocked\\
Corpus match & measured after generation & measured after generation\\
Grammar update from output & never automatic & never automatic\\
Provenance & prompt/output audit & proposal + licenses + realization\\
\bottomrule
\end{tabular}
\end{table}

\subsection{Research questions}

The experiment asks three related questions: (i) what kinds of divergence
appear under direct LLM surface generation; (ii) whether an LLM can preserve
the requested abstract structural specification when surface realization is
delegated to reviewed grammar; and (iii) whether any proposal in the frozen
task set lacks a required license and is therefore blocked. In v1, all 28
proposals preserve the requested specification and all 28 are realizable.
Accordingly, no blocking event occurs in this deliberately licensed task set;
the experiment does not estimate a general blocking rate for arbitrary Bororo
requests.

\subsection{Experimental controls}

To make the two conditions comparable, task specifications, target lexical
items, person values, and constructions were drawn from the same frozen
inventory. Prompts and model settings were frozen and archived with the
experiment. The direct condition did not receive additional examples unavailable
to the proposal condition. Raw outputs were retained, and classifications that
depend on linguistic analysis remain auditable against the reviewed grammar.

The v1 grammar and task inventory were frozen before the LLM outputs were
collected. Inspection of those outputs did not modify the grammar used to score
the same experiment. Linguistic observations arising during audit are therefore
eligible only for a later review version. This prevents evaluation data from
silently changing the grammatical licenses against which they are evaluated.

\subsection{Evaluation categories}

Outputs are classified independently along several dimensions rather than
collapsed into a single accuracy score. Crucially, morphological well-formedness
is separated from task adequacy. A grammatical Bororo form may instantiate a
different construction from the one requested. For example, a form without the
indicative marker need not be morphologically ill formed merely because the
experimental task requested indicative predication.

The audit therefore distinguishes \emph{morphological violation},
\emph{construction mismatch}, position-specific \emph{S-, A-, and O-person
mismatch}, and \emph{frame violation}. It additionally records
\emph{constructional overgeneralization}, \emph{complementary-distribution
violation}, and \emph{unsupported but plausible} analyses. Unresolved cases
remain unresolved rather than being forced into an error category.

Corpus status is reported separately. A candidate may be sentence-attested,
predicate-form-attested, or unattested in the supplied corpus. These labels
describe documentary evidence and do not alter its grammatical-review status.

\subsection{Results of the frozen v1 comparison}

The frozen experiment contains 28 tasks. In the direct condition, only one
output (1/28; 3.6\%) is an exact surface match to the deterministic realization.
This number is a string-level comparison and is \emph{not} interpreted as a
3.6\% grammaticality or accuracy score. Human audit currently marks 12 outputs
as reviewed and 16 as pending further linguistic analysis. Twenty-five of the
28 direct outputs show a construction mismatch relative to the requested task,
whereas three do not.

The multidimensional audit illustrates why exact nonmatch cannot be equated
with ungrammaticality. Morphological violation is established for three outputs,
rejected for nine, and remains unresolved for sixteen. For S-person realization,
nine outputs are mismatches, ten are not, and three remain unresolved among the
22 applicable S tasks. All six divalent tasks show an A-person mismatch. For O,
two of the six divalent tasks match, two mismatch, and two remain unresolved.
Six outputs show a coding-frame violation, eighteen do not, and four remain
unresolved. No complementary-distribution violation was identified in the 28
outputs. Constructional overgeneralization and unsupported-but-plausible status
remain unresolved in fourteen and sixteen cases, respectively.

In the evidence-bounded condition, all 28 abstract proposals preserve the
requested task specification exactly: lemma, construction, S, A, O, and oblique
fields show no proposal--task mismatch. The frozen deterministic system then
realizes all 28 proposals, with no blocked request. Thus the observed v1 chain
is 28 tasks $\rightarrow$ 28 structurally aligned proposals $\rightarrow$ 28
licensed deterministic realizations.

\begin{table}[ht]
\centering
\caption{Frozen v1 LLM comparison. Audit categories are reported independently.}
\label{tab:v1llmresults}
\begin{tabular}{lr}
\toprule
Measure & Result\\
\midrule
Tasks & 28\\
Direct exact surface matches & 1/28 (3.6\%)\\
Direct construction mismatches & 25/28\\
Direct morphological violation & 3 yes; 9 no; 16 unresolved\\
Direct S-person mismatch (22 applicable) & 9 yes; 10 no; 3 unresolved\\
Direct A-person mismatch (6 applicable) & 6 yes\\
Direct O-person mismatch (6 applicable) & 2 yes; 2 no; 2 unresolved\\
Direct frame violation & 6 yes; 18 no; 4 unresolved\\
Direct complementary-distribution violation & 0/28\\
Condition-B proposals aligned to task & 28/28\\
Condition-B proposals deterministically realized & 28/28\\
\bottomrule
\end{tabular}
\end{table}

These results test task preservation and evidence-bounded realization, not
native-speaker acceptability. In particular, a deterministic realization means
that the frozen reviewed system licensed the requested structural cell; it is
not an independent grammaticality judgment. Conversely, a direct output that
differs from the controlled realization may still be a well-formed Bororo
expression in another construction.

\n\section{Evaluation design and limitations}
\label{sec:limits}

The baseline is intentionally small. It measures the behavior of an
evidence-bound generation architecture, not general coverage of Bororo.
Predicate-form occurrence is morphological evidence only; in particular, for a
divalent predicate it does not resolve the A argument of the generated
structural cell. Conversely, a generated form absent from the supplied corpus
is not thereby classified as ungrammatical.

The same caution applies to the LLM comparison. Exact agreement with the
controlled surface is a reproducible string-level measure, but disagreement is
not itself evidence of ungrammaticality. Direct outputs are therefore audited
along independent dimensions, and an unresolved analysis is retained as such.
The current human audit is incomplete by design where the frozen evidence does
not justify a stronger classification.

The reported numbers should therefore be read as coverage, task-preservation,
and evidence-control measures rather than as a single accuracy or
grammaticality score. Broader evaluation requires additional reviewed lexical
items and constructions and, for claims about acceptability beyond the reviewed
analysis, independent speaker judgment.

\section{Discussion}

The two experiments expose different consequences of separating evidence from
licensing. The controlled baseline shows that partial reviewed knowledge can
support productive recombination without collapsing corpus attestation into a
grammar: all 28 requested cells were realized although no complete generated
sentence was an exact corpus sentence. At the same time, most baseline records
depend on individually reviewed cells rather than paradigm-wide extrapolation.
The architecture therefore represents incompleteness explicitly instead of
silently completing paradigms by analogy.

The LLM comparison adds a second result. Direct surface generation and abstract
structural proposal behave differently under the same frozen task inventory.
Only one direct output is an exact match to the controlled realization, but this
fact cannot be interpreted as 27 grammatical errors. The human audit shows why:
some direct outputs are potentially well-formed expressions that fail the
requested construction or person specification, while the internal analysis of
others remains unresolved. In particular, the distinction between
morphological well-formedness and task adequacy is empirically necessary rather
than merely terminological.

Condition B isolates a narrower capability of the LLM. All 28 abstract
proposals preserve the requested structural specification, and all 28 can then
be realized from the grammar frozen before the model outputs were collected.
The result does not show that the deterministic system has complete Bororo
coverage, nor does it constitute independent native-speaker validation.
Instead, it shows that, for this reviewed test set, probabilistic structural
proposal can be separated from surface morphological realization without loss
of the requested task information. The model need not be entrusted with the
morphological generalizations that the experiment is intended to control.

The asymmetry is important for low-resource work. A direct model output can be
fluent-looking while leaving unclear whether a divergence reflects
morphological ill-formedness, a different grammatical construction, incorrect
participant indexing, or an analysis not yet supported by the available
description. An evidence-bounded system can instead preserve these states as
distinct outcomes. Blocking and unresolved status are consequently informative
behaviors, not merely failures to generate.

The present comparison remains deliberately small. It contains five target
lexemes and 28 reviewed tasks, and 16 direct outputs still require further
linguistic analysis in at least part of the audit. The 28/28 realization result
for Condition B is also conditioned on tasks sampled from the frozen reviewed
inventory; it should not be generalized to arbitrary Bororo structures. A
larger experiment should add lexical items and construction types while
preserving the same pre-registration principle: grammar and task inventory are
frozen before model outputs are inspected.

Finally, the architecture supports incremental revision without contaminating
an existing evaluation. New corpus observations can enter a review queue
without changing the grammar. Human analysis can subsequently add an exact
cell, a construction-specific realization, or, where independently justified,
a complete lexical class to a later grammar version. Experimental discoveries
therefore remain evidence for future analysis rather than retroactive licenses
for the outputs being evaluated.

\section{Conclusion}

This study argues for an evidential boundary between documentary data,
linguistic analysis, and computational generation in low-resource language
technology. In the Bororo case study, corpus occurrence is observational
evidence rather than an automatic productive rule, and human-reviewed
grammatical knowledge is represented explicitly enough to license, block, or
leave a request unresolved.

The frozen baseline demonstrates controlled recombination across 28 reviewed
structural cells without exact sentence retrieval from the corpus. The LLM
experiment then shows a complementary result: direct surface generation rarely
matches the controlled realization exactly and requires multidimensional human
audit, whereas all 28 abstract proposals preserve the requested structure and
can be realized by the frozen evidence-bounded system. These findings concern
task preservation and reviewed licensing, not a general grammaticality score or
a claim of complete coverage.

The broader methodological implication is that an LLM need not function as the
grammar of a low-resource language in order to contribute to generation. It can
operate as a proposal mechanism upstream of an explicit, revisable, and
auditable grammatical layer. This arrangement keeps uncertainty visible,
prevents experimental outputs from becoming their own evidence, and provides a
reproducible path for extending computational resources as linguistic analysis
progresses.

\bibliographystyle{plain}
\bibliography{references}

\end{document}
